% 金星（综述）
% license CCBYSA3
% type Wiki

本文根据 CC-BY-SA 协议转载翻译自维基百科\href{https://en.wikipedia.org/wiki/Venus}{相关文章}。

金星是距离太阳第二近的行星。由于其轨道最接近地球，且两者同为类地岩石行星，并在大小、质量和表面引力上几乎完全一致，因此金星常被称为地球的“孪生”或“姐妹”行星。然而，金星在许多方面与地球显著不同，尤其是它没有液态水，其大气层远比太阳系中其他任何岩石天体都更加厚重和致密。金星大气主要由二氧化碳组成，并覆盖着一层浓密的硫酸云层，环绕整个行星。在平均地表高度处，大气温度达到 737\,K（464\,°C；867\,°F），压力为地球海平面的 92 倍，使其最低层大气呈现超临界流体状态。从地球观测时，金星呈现为类似恒星的亮点，比天空中任何其他天然光源都更明亮\(^\text{[22][23]}\)，并且作为一颗内行星，它总是出现在接近太阳的方向，要么作为最明亮的“启明星”，要么作为“长庚星”。

金星与地球的轨道使两者在约 1.6 年的会合周期中彼此接近。在此过程中，金星会比其他任何行星更接近地球；相较之下，由于水星的轨道更靠近太阳，因此它在自身轨道路径中整体上比其他行星更接近地球。在从地球出发的行星际空间飞行中，金星常被用作引力助推的中转点，从而提供更快速且更经济的航程。金星没有卫星，并以非常缓慢的逆行方式绕其自转轴旋转，这被认为是太阳潮汐锁定效应与金星巨大大气层的受热差异相互竞争的结果。因此，金星的一昼夜长达 116.75 个地球日，而其一个太阳年为 224.7 个地球日。

金星具有微弱的磁层；由于缺乏内部发电机，其磁场是由太阳风与大气相互作用所诱导产生的。在内部结构上，金星具有核心、地幔和地壳。内部热量通过活跃的火山活动逸出\(^\text{[24][25]}\)，从而导致地表更新，而不是板块构造。金星在其早期历史中可能曾拥有液态地表水及宜居环境\(^\text{[26][27]}\)，但由于失控的温室效应，所有水分最终蒸发，使金星演变为当前的状态\(^\text{[28][29][30]}\)。在云层高度处的大气条件是整个太阳系中与地球最为相似的，这些条件被认为可能有利于金星生命的存在，并在 2020 年发现了潜在的生物标志物，从而推动了新的研究与探测任务。

在人类历史上，全世界的人们都曾观测过金星，并使其在众多文化中具有特殊的重要性。随着望远镜的出现，金星的相位变得可辨识，并在 1613 年被呈现为推翻当时占主导地位的地心模型、支持日心模型的决定性证据。1961 年，“金星 1 号”首次造访金星，飞掠行星并实现了人类第一次行星际飞行。1962 年，第二次行星际任务“水手 2 号”带回了首批来自金星的数据。1967 年，首个行星际撞击器“金星 4 号”抵达金星，随后在 1970 年由着陆器“金星 7 号”成功登陆。截至 2025 年，“太阳轨道飞行器”正前往于 2026 年飞掠金星，而下一项计划发射前往金星的任务是“金星生命探测器”，预定同样于 2026 年发射。
\subsection{物理特征}
\begin{figure}[ht]
\centering
\includegraphics[width=8cm]{./figures/23fb842f59355847.png}
\caption{金星（从左数第二，伪彩色显示）按比例与太阳系内侧具有行星质量的天体一起展示，排列顺序依其轨道自太阳向外依次排列（从左到右：水星、金星、地球、月球、火星和谷神星）。} \label{fig_Venus_1}
\end{figure}
\begin{figure}[ht]
\centering
\includegraphics[width=8cm]{./figures/f163eb810133acdf.png}
\caption{金星在不同波长下的成像} \label{fig_Venus_2}
\end{figure}
金星是太阳系中四颗类地行星之一，这意味着它与地球一样属于岩石天体。它在大小和质量上与地球相似，因此常被描述为地球的“姐妹”或“孪生”行星\(^\text{[31]}\)。由于其自转缓慢，金星的形状非常接近球形\(^\text{[32]}\)。其直径为 12,103.6\,km（7,520.8\,mi），仅比地球小 638.4\,km（396.7\,mi）；其质量为地球的 81.5\%，使金星成为太阳系中第三小的行星。金星表面的环境与地球完全不同，因为其致密大气中 96.5\% 为二氧化碳，导致强烈的温室效应，其余约 3.5\% 为氮气\(^\text{[33]}\)。金星表面气压为 9.3\,MPa（93\,bar），平均表面温度为 737\,K（464\,°C；867\,°F），高于两种主要组分的临界点，使其地表大气成为一种超临界流体，主要由超临界二氧化碳及部分超临界氮组成。
\subsubsection{自然历史}
\textbf{形成}

包括金星在内的类地岩石行星被认为经历了五个阶段形成：尘埃沉降、微行星形成、行星胚胎阶段、巨撞阶段，以及最终的大气形成阶段。来自金星的有限测量数据使其形成时间线难以进行更详细的分析\(^\text{[34]}\)。

\textbf{未来}

金星预计将在七至八十亿年后当太阳演变为红巨星时，与水星一起被毁灭，也有可能包括地球和月球\(^\text{[35]}\)。
\subsubsection{地理}
\begin{figure}[ht]
\centering
\includegraphics[width=8cm]{./figures/884bdec4a9f38b8d.png}
\caption{彩色编码的高程地图，显示金星上呈黄色的高地“大陆”以及次要地形特征} \label{fig_Venus_3}
\end{figure}
\begin{figure}[ht]
\centering
\includegraphics[width=8cm]{./figures/4ef3c5c0b6c896cf.png}
\caption{金星表面雷达数据的球形视图，突出显示地表特征（1989 年，麦哲伦号）。这些颜色并不代表地表的真实外观[36]。} \label{fig_Venus_4}
\end{figure}
在 20 世纪探测器揭示部分秘密之前，金星表面一直是推测的对象。1975 年和 1982 年的“金星号”着陆器传回了由沉积物和相对棱角分明的岩石覆盖的地表图像\(^\text{[37]}\)。1990–91 年间，“麦哲伦号”对其表面进行了详细测绘。有大量火山活动的证据，并且大气中二氧化硫的变化可能表明存在活火山\(^\text{[38][39]}\)。

大约 80\% 的金星表面被光滑的火山平原覆盖，其中 70\% 为带有皱褶脊的平原，10\% 为光滑或叶状平原\(^\text{[40]}\)。其余的地表由两个高地“大陆”构成，一个位于行星的北半球，另一个位于赤道以南。北部高地称为伊什塔尔高原（Ishtar Terra），得名于巴比伦的爱神伊什塔尔，其面积约与澳大利亚相当。麦克斯韦山脉位于伊什塔尔高原，其最高峰斯卡迪山（Skadi Mons）是金星上的最高点，海拔比金星平均表面高度高 11\,km（7\,mi）\(^\text{[41]}\)。南部高地称为阿芙洛狄忒高原（Aphrodite Terra），得名于希腊神话中的爱神，是两大高地中较大的一个，面积约与南美洲相当。该区域的大部分被裂缝和断层网络覆盖\(^\text{[42]}\)。

近期有关于金星上岩浆流动的证据（2024 年）\(^\text{[43]}\)，例如在盾状火山 Sif Mons 的熔岩流以及在平坦平原 Niobe Planitia 上的熔岩流\(^\text{[24]}\)。地表可见破火山口。金星的撞击坑数量很少，表明其表面相对年轻，约为 3 亿至 6 亿年\(^\text{[44][46]}\)。金星除了拥有岩石行星上常见的撞击坑、山脉和山谷外，还具有一些独特的地表特征。其中包括称为“farra”的平顶火山构造，其形状类似煎饼，直径从 20 至 50\,km（12 至 31\,mi）不等，高度在 100 至 1,000\,m（330 至 3,280\,ft）之间；称为“novae”的放射状、星形裂缝系统；具有放射状和同心裂隙、类似蜘蛛网的构造，被称为“arachnoids”；以及“coronae”，即被洼地包围的环状裂缝结构。这些特征均具有火山成因\(^\text{[46]}\)。
\begin{figure}[ht]
\centering
\includegraphics[width=8cm]{./figures/5421d90f38272eab.png}
\caption{彩色化影像（“金星 9 号”，1975 年），金星天空在地表呈橙黄色，这是由于瑞利散射或低层大气中的蓝色吸收体所致，而在更高的高度则呈白色[47][48]；同时其地表为类似玄武岩的深灰色，可能因氧化而呈红色[36]。} \label{fig_Venus_5}
\end{figure}
大多数金星表面特征以历史和神话中的女性命名\(^\text{[49]}\)。例外包括麦克斯韦山脉，以詹姆斯·克拉克·麦克斯韦命名，以及阿尔法高地区（Alpha Regio）、贝塔高地区（Beta Regio）和奥夫达高地区（Ovda Regio）。这后三个特征是在国际天文学联合会（负责行星命名的机构）采用现行系统之前命名的\(^\text{[50]}\)。

金星地形特征的经度以其本初子午线为基准表达。最初的本初子午线通过位于阿尔法高地区南部、椭圆形地形“夏娃”（Eve）中央的雷达亮点\(^\text{[51]}\)。在“金星号”任务完成后，本初子午线被重新定义为通过塞德娜平原（Sedna Planitia）上阿里阿德涅（Ariadne）陨石坑中央峰的位置\(^\text{[9][52]}\)。

地层学上最古老的镶嵌地形（tessera terrains）在“金星快车号”和“麦哲伦号”测量中显示出相对于周围玄武质平原更低的热辐射率，这表明其具有不同的、可能更富长英质（felsic）的矿物组合\(^\text{[29][53]}\)。生成大量长英质地壳的机制通常需要存在海洋和板块构造，这意味着早期金星上可能存在宜居条件，并在某一时期拥有大量水体\(^\text{[54]}\)。然而，镶嵌地形的具体性质仍不确定\(^\text{[55]}\)。

2023 年的研究首次提出金星在古代可能存在板块构造，因此可能拥有更宜居的环境，甚至可能具备维持生命的能力\(^\text{[26][27]}\)。金星因此成为研究类地行星形成及其宜居性的重要案例。

\textbf{火山活动}
\begin{figure}[ht]
\centering
\includegraphics[width=8cm]{./figures/96e9067c358877bf.png}
\caption{金星埃斯特拉区域内两个煎饼穹丘的雷达拼接图——二者均宽 65\,km（40\,mi），高度不足 1\,km（0.62\,mi）} \label{fig_Venus_6}
\end{figure}
金星表面的很大一部分似乎是由火山活动塑造的。金星上的火山数量是地球的数倍，并拥有 167 座直径超过 100\,km（60\,mi）的大型火山。地球上唯一与其规模相当的火山构造是夏威夷大岛\(^\text{[46]:154}\)。金星上已识别并绘制地图的火山数量超过 85,000 座\(^\text{[56][57]}\)。这并不是因为金星的火山活动比地球更活跃，而是因为其地壳更古老，且没有经历地球上活跃的侵蚀过程。地球的海洋地壳在构造板块边界通过俯冲不断被循环，其平均年龄约为 1 亿年\(^\text{[58]}\)，而金星表面估计为 3 亿至 6 亿年\(^\text{[44][46]}\)。

多项证据指向金星上持续存在火山活动。上层大气中的二氧化硫浓度在 1978 年至 1986 年间下降了 10 倍，在 2006 年急剧上升，然后再次下降 10 倍\(^\text{[59]}\)。这可能意味着多次大型火山喷发提高了其含量\(^\text{[60][61]}\)。有人提出金星的闪电（见下文）可能源自火山活动（即火山闪电）。2020 年 1 月，天文学家报告发现证据表明金星目前存在火山活动，具体而言，他们检测到橄榄石这种火山产物，而该物质在金星表面会迅速风化\(^\text{[62][63]}\)。

这一巨大的火山活动由炽热的内部提供能量，根据模型，它可以由行星早期阶段的高能碰撞以及类似地球的放射性衰变来解释。撞击的速度会明显高于地球，一方面因为金星更靠近太阳而运动更快，另一方面因为高偏心率天体与金星相撞时具有更高速度\(^\text{[64]}\)。

在 2008 年和 2009 年，“金星快车号”首次观测到正在进行火山活动的直接证据，其形式是在裂谷区加尼斯峡谷（Ganis Chasma）内的四个短暂出现的局部红外热点\(^\text{[note 1]}\)，该区域靠近盾状火山玛特山（Maat Mons）\(^\text{[65]}\)。其中三个热点在连续轨道中被重复观测到。这些热点被认为代表新近喷发的熔岩\(^\text{[66][67]}\)。实际温度未知，因为热点的面积无法测量，但推测其温度可能在 800–1,100\,K（527–827\,°C；980–1,520\,°F）范围内，而正常温度约为 740\,K（467\,°C；872\,°F）\(^\text{[68]}\)。2023 年，科学家重新分析了“麦哲伦号”拍摄的玛特山区域地形图像。通过计算机模拟，他们确认该区域在 8 个月间发生了地形变化，并得出其原因是活跃火山活动\(^\text{[69]}\)。

\textbf{陨石坑}
\begin{figure}[ht]
\centering
\includegraphics[width=8cm]{./figures/d8116d164677ab95.png}
\caption{金星表面的撞击坑（伪彩色，由雷达数据重建的三维投影影像）} \label{fig_Venus_7}
\end{figure}
金星上有将近一千个撞击坑，均匀分布在其表面。在其他具有撞击坑的天体上，例如地球和月球，撞击坑呈现出不同程度的退化状态。月球上的退化由后续撞击造成，而地球上的退化则由风化和降雨侵蚀导致。在金星上，约 85\% 的撞击坑保持完好状态。撞击坑的数量及其良好的保存状况表明，该行星曾在 3 亿至 6 亿年前经历过一次全球性地表重塑事件\(^\text{[44][45]}\)，随后火山活动逐渐衰减\(^\text{[70]}\)。地球的地壳处于持续运动中，而金星被认为无法维持此类过程。在没有板块构造来散发地幔热量的情况下，金星经历一种循环过程，即地幔温度不断升高，直到达到削弱地壳的临界水平。随后，在大约 1 亿年的时间内，发生大规模的俯冲，将地壳完全循环\(^\text{[46]}\)。

金星的撞击坑直径范围从 3 到 280\,km（2 到 174\,mi）。没有直径小于 3\,km 的撞击坑，这是由于致密大气对进入物体的影响所致。动能低于某一阈值的物体会被大气大幅减速，以至于无法形成撞击坑\(^\text{[71]}\)。直径小于 50\,m（160\,ft）的入射物体将在到达地面之前于大气层中解体并焚毁\(^\text{[72]}\)。
\subsubsection{内部结构}
\begin{figure}[ht]
\centering
\includegraphics[width=8cm]{./figures/dbfa21d636bc841c.png}
\caption{金星的分化结构} \label{fig_Venus_8}
\end{figure}
由于缺乏反射地震学的数据或其转动惯量的测量结果，人们对金星内部结构和地球化学性质的直接信息非常有限\(^\text{[73]}\)。金星与地球在大小和密度上的相似性表明，它们可能拥有相似的内部结构：核心、地幔和地壳。与地球相似，金星的核心很可能至少部分为液态，因为这两颗行星的冷却速率大致相同\(^\text{[74]}\)，尽管完全固态核心也不能被排除\(^\text{[75]}\)。金星体积略小意味着其深部内部的压强比地球低 24\%\(^\text{[76]}\)。根据行星模型预测的转动惯量数值，金星核心半径估计为 2,900–3,450\,km\(^\text{[75]}\)。基于 2006 至 2020 年之间测量的自转轴进动速率所推算的转动惯量，目前的估计为 3,500\,km\(^\text{[13][77]}\)。

金星地壳平均估计厚度为 40\,km，最大可达 65\,km\(^\text{[78]}\)。

两颗行星的主要差异在于金星缺乏板块构造的证据，可能原因是其地壳太坚固，在没有水减少黏滞性的情况下无法发生俯冲。这导致行星热量损失较少，阻碍其冷却，并为其缺乏内部磁场提供了可能的解释\(^\text{[79]}\)。金星可能改以周期性的全球地表重塑事件释放内部热量\(^\text{[44]}\)。
\subsubsection{磁场与核心}
\begin{figure}[ht]
\centering
\includegraphics[width=8cm]{./figures/81a077950d438b05.png}
\caption{} \label{fig_Venus_9}
\end{figure}
1967 年，“金星 4 号”发现金星的磁场比地球弱得多。该磁场是由电离层与太阳风之间的相互作用所诱导产生的\(^\text{[80][81][page needed]}\)，而不是像地球那样由核心内部的发电机产生。金星这种微弱的诱导磁层对大气抵御太阳和宇宙辐射的保护几乎可以忽略不计。

金星缺乏本征磁场令科学家感到意外，因为金星在大小上与地球相似，原本应在其核心中存在发电机。发电机需要三个条件：导电液体、自转和对流。金星的核心被认为具有电导性，而尽管其自转通常被认为过于缓慢，模拟表明其仍足以产生发电机\(^\text{[82][83]}\)。这意味着缺乏发电机的原因在于金星核心内部缺乏对流。在地球上，对流发生在核心的液态外层中，因为液态层底部的温度远高于顶部。在金星上，一次全球性的地表重塑事件可能终止了板块构造，导致通过地壳的热通量减小。这种隔热效应会使地幔温度升高，从而减少核心的热通量。因此，无法提供用于驱动磁场的内部地球发电机，核心释放的热量则重新加热地壳\(^\text{[84]}\)。

一种可能性是金星没有固态内核\(^\text{[85]}\)，或者其核心没有在冷却，使得整个液态核心的温度大致相同。另一种可能性是其核心已经完全固化。核心的状态高度依赖于硫的浓度，而当前尚不清楚\(^\text{[84]}\)。

另一种可能性是金星缺乏像地球“形成月球”事件那样的大型撞击，使金星核心在增量形成过程中产生分层结构，而没有足够的力量触发或维持对流，从而无法形成“地质发电机”\(^\text{[86]}\)。

金星周围微弱的磁层意味着太阳风会直接与其外层大气相互作用。在这里，水分子因紫外线辐解而生成氢和氧离子。随后，太阳风提供能量，使其中一些离子获得足够速度逃离金星的引力场。该侵蚀过程导致轻质的氢、氦和氧离子持续流失，而高质量分子如二氧化碳更可能被保留。太阳风导致的大气侵蚀可能造成金星在形成后的前十亿年内失去了大部分水\(^\text{[87]}\)。然而，该行星在最初的 20–30 亿年内可能保有发电机，因此水的流失也可能发生得更晚\(^\text{[88]}\)。这种侵蚀使得大气中高质量氘与低质量氢的比值相比太阳系其他区域增加了 100 倍\(^\text{[89]}\)。
\subsection{大气与气候}
\begin{figure}[ht]
\centering
\includegraphics[width=8cm]{./figures/8e5406aad8f1ba0c.png}
\caption{通过紫外成像显现的金星大气云层结构} \label{fig_Venus_10}
\end{figure}
金星拥有致密的大气层，由 96.5\% 的二氧化碳、3.5\% 的氮气组成——两者在金星地表均以超临界流体形式存在，其密度为水的 6.5\%\(^\text{[90]}\)——并含有微量其他气体，包括二氧化硫\(^\text{[91]}\)。其大气质量为地球的 92 倍，而地表压力约为地球的 93 倍，相当于地球海洋表面下近 1\,km（5⁄8\,mi）处的压力。地表大气密度为 65\,kg/m³（4.1\,lb/cu\,ft），为水的 6.5\%\(^\text{[90]}\)，或为地球海平面 293\,K（20\,°C；68\,°F）空气密度的 50 倍。富含 CO₂ 的大气产生了太阳系中最强的温室效应，使其表面温度至少达到 735\,K（462\,°C；864\,°F）\(^\text{[92][93]}\)。这使得金星表面比水星更热\(^\text{[94]}\)，尽管金星与太阳的距离几乎是水星的两倍，因此其接收的太阳辐照度仅为水星的四分之一，即 2,600\,W/m²（约为地球的两倍）\(^\text{[5]}\)。由于其失控温室效应，金星被卡尔·萨根等科学家视为与地球气候变化相关的警示和研究对象\(^\text{[95]}\)。因此金星被称为温室行星\(^\text{[96]}\)，或处于温室炼狱中的行星\(^\text{[97]}\)。

金星大气中富含原初稀有气体，与地球大气相比更为丰富\(^\text{[98]}\)。这种富集表明金星在演化早期便与地球发生分歧。一个异常大的彗星撞击\(^\text{[99]}\)或从太阳星云吸积更大质量的原始大气\(^\text{[100]}\)被提出用于解释这种富集。然而，大气中放射性氩–40 含量较低，这种元素可视为地幔脱气的替代指标，这表明主要岩浆作用在早期便已终止\(^\text{[101][102]}\)。

研究表明，在数十亿年前，金星的大气可能与早期地球的大气非常相似，并且其地表可能存在大量液态水\(^\text{[103][104][105]}\)。经过约 6 亿年至数十亿年的时间\(^\text{[106]}\)，太阳光度上升及可能发生的大规模火山重塑导致原始水体蒸发\(^\text{[107]}\)。当大气中的温室气体（包括水）达到临界值后，便触发了失控温室效应\(^\text{[108]}\)。尽管金星地表环境已不再适合任何可能在此事件前形成的类地生命，但人们推测金星表面上方 50\,km（30\,mi）的高层云层中可能存在生命，因为该高度的环境是整个太阳系中最类似地球的\(^\text{[109]}\)，温度介于 303 至 353\,K（30 至 80\,°C；86 至 176\,°F）之间，气压和辐射水平与地球表面相当，但大气成分为酸性云和二氧化碳\(^\text{[110][111][112]}\)。更具体而言，在 48 至 59\,km 的高度范围内，温度和辐射条件适宜生命生存；在更低的高度水会蒸发，而在更高的高度紫外辐射将过强\(^\text{[113][114]}\)。2020 年 9 月，有研究报告在金星大气中探测到磷化氢的吸收线，而该物质在已知条件下无法通过非生物途径生成，因此引发了关于金星大气中可能存在现存生命的猜测\(^\text{[115][116]}\)。后续研究将这一光谱信号解释为二氧化硫\(^\text{[117]}\)，或发现实际上不存在该吸收线\(^\text{[118][119]}\)。
\begin{figure}[ht]
\centering
\includegraphics[width=8cm]{./figures/b77f963239de7576.png}
\caption{按高度绘制的大气剖面图（左侧刻度）：云层、温度变化（VIRA 粗线及底部刻度）、压力变化（右侧刻度）以及风速（PV 虚线及顶部刻度）} \label{fig_Venus_11}
\end{figure}
热惯性以及下层大气中由风输送热量的作用意味着，尽管金星自转缓慢，其面对太阳与背向太阳的半球之间的地表温度并没有显著差异。地表的风速很慢，仅为每小时数公里，但由于地表大气密度极高，它们对障碍物施加的力非常大，并能够在地表搬运尘埃和小石块。即便不考虑其高温、高压和缺氧环境，这种阻力本身就足以使人在地表难以行走\(^\text{[120]}\)。

在致密的 CO₂ 层之上，距地表 45 至 70\,km 的高度存在厚厚的云层\(^\text{[121]}\)，主要由硫酸组成，其形成过程是二氧化硫分子在紫外辐射催化下与水发生反应\(^\text{[122]}\)，从而生成硫酸水合物\(^\text{[123]}\)。此外，云层中含有约 1\% 的三氯化铁\(^\text{[124][125]}\)。云粒子的其他可能成分包括硫酸铁、氯化铝和五氧化二磷。不同高度的云层具有不同的成分和粒径分布\(^\text{[124]}\)。这些云层像地球的厚云层一样，将落在其上的阳光约 70\% 反射回太空\(^\text{[1265]}\)，并由于覆盖整个行星，使得金星表面无法被可见光直接观测。永久性的云层意味着尽管金星比地球更接近太阳，其地表接收到的阳光却更少，只有 10\% 的入射阳光能到达地表\(^\text{[127]}\)，使得其地表白天平均照度约为 14,000 勒克斯，与地球“阴天白昼”的照度相当\(^\text{[128]}\)。

金星大气的自转速度远快于其固体本体的自转，这一现象被称为大气超自转\(^\text{[129]}\)。这导致云顶处出现强劲的 300\,km/h（185\,mph）风，它们绕行星一周仅需约 4 天，相当于行星自转速度的 60 倍\(^\text{[130]}\)，而地球上最强的风速仅为其自转速度的 10–20\%。

尽管金星在可见光下看起来缺乏表面特征，但在紫外波段却存在条带或条纹，其成因尚未确定。紫外吸收可能来源于一种由氧与硫组成的化合物 OSSO，该分子中两个硫原子之间具有双键，并存在“顺式”和“反式”两种构型，或可能来源于 S₂ 至 S₈ 的多硫化合物\(^\text{[131]}\)。
\begin{figure}[ht]
\centering
\includegraphics[width=8cm]{./figures/d0e3e9800a1218c4.png}
\caption{} \label{fig_Venus_12}
\end{figure}
金星的地表实际上是等温的；它不仅在两个半球之间保持恒定温度，而且在赤道与两极之间也保持恒定温度\(^\text{[5][131]}\)。金星的自转轴倾角极小——不到 3°，相比之下地球为 23°——这也使季节性温度变化极小\(^\text{[134]}\)。海拔是影响金星温度为数不多的因素之一。金星最高点——麦克斯韦山脉中的斯卡迪山——因此是金星上最冷的地点，温度约为 655\,K（380\,°C；715\,°F），气压约为 4.5\,MPa（45\,bar）\(^\text{[132][135]}\)。1995 年，“麦哲伦号”探测器在金星最高山峰顶部成像到一种高反射物质，被称为“金星雪”，与地球上的雪极为相似。这种物质可能通过类似于雪的过程形成，但温度高得多。由于其挥发性过高而无法在地表凝结，它以气态上升至更高、更冷的高度，在那里可能发生沉降。该物质的成分尚未确定，但推测范围从碲单质到硫化铅（方铅矿）不等\(^\text{[136]}\)。

尽管金星没有季节变化，天文学家在 2019 年发现其大气对阳光吸收存在周期性变化，可能由悬浮在高层云中的不透明吸收粒子引起。这种变化导致金星纬向风速出现观测到的变化，并似乎与太阳 11 年黑子周期同步\(^\text{[137]}\)。

金星大气中是否存在闪电长期以来存在争议\(^\text{[138]}\)，自从苏联“金星号”探测器首次探测到疑似闪电脉冲后\(^\text{[139][140][141]}\)。在 2006–07 年，“金星快车号”清晰探测到了哨声模式波，这是闪电的特征信号。其间歇性出现表明其与天气活动相关。根据这些测量，其闪电发生率至少为地球的一半\(^\text{[142]}\)，然而其他仪器完全未能探测到闪电\(^\text{[139]}\)。闪电的起源仍不明确，但可能来自云层或金星火山。

2007 年，“金星快车号”发现金星南极存在一个巨大的双极大气涡旋\(^\text{[143][144]}\)。“金星快车号”在 2011 年发现金星高层大气中存在臭氧层\(^\text{[145]}\)。2013 年，欧洲航天局科学家报告称金星的电离层向外流动，其方式类似于“在类似条件下彗星电离尾的外泄”\(^\text{[146][147]}\)。

在 2015 年 12 月，以及较弱的 2016 年 4 月和 5 月，参与日本“晓号”任务的研究人员在金星大气中观测到弧形结构。这被认为是太阳系中最大、可能也是首个直接证据确认的驻波重力波\(^\text{[148][149][150]}\)。

金星地表大气的颜色与声音\(^\text{[151]}\)已被记录：天空在地表呈橙黄色，而在更高的高度则呈白色\(^\text{[48]}\)。
\subsection{轨道与自转}
金星以约 0.72\,AU（1.08 亿\,km；6700 万\,mi） 的平均距离绕太阳运行，每 224.7 天完成一周。它在 7.998 年内完成 13 圈公转，因此在天空中的位置几乎每八年重复一次。尽管所有行星轨道都是椭圆形的，但金星的轨道目前最接近圆形，其偏心率小于 0.01\(^\text{[5]}\)。对早期太阳系轨道动力学的模拟显示，金星轨道的偏心率在过去可能大得多，高达 0.31，并可能影响其早期气候演化\(^\text{[152]}\)。
\begin{figure}[ht]
\centering
\includegraphics[width=8cm]{./figures/cefd4e3932d4f154.png}
\caption{金星及其自转相对于公转的关系} \label{fig_Venus_13}
\end{figure}
金星具有逆行自转，这意味着它绕自身轴旋转的方向与包括地球在内的大多数行星不同，为顺时针方向，与其绕太阳的逆时针公转方向相反。因此金星的恒星日为 243 个地球日，比其公转周期 224.7 个地球日更长。如果金星被潮汐锁定在太阳上，它将始终以同一面朝向太阳，其恒星日将为 224.7 天。然而，由于金星大气质量巨大且接近太阳，大气的差异加热使金星具有微弱的逆行自转。由于同样的原因，其昼长也可波动长达 20 分钟\(^\text{[153][154]}\)。利用“麦哲伦号”探测器数据在 500 天期间测量得到的金星自转周期小于在“麦哲伦号”和“金星快车号”访期间隔 16 年中测得的自转周期，两者相差约 6.5 分钟\(^\text{[155]}\)。由于逆行自转，金星的太阳日显著短于其恒星日，仅为 116.75 个地球日\(^\text{[12]}\)。一个金星年约相当于 1.92 个金星太阳日\(^\text{[156]}\)。对于位于金星地表的观察者而言，太阳将从西方升起并在东方落下\(^\text{[156]}\)，尽管金星不透明的云层使地表无法看到太阳\(^\text{[157]}\)。

金星可能在太阳星云中形成时具有不同的自转周期和自转轴倾角，并由于行星摄动和潮汐效应对其致密大气产生的混沌自转变化而达到当前状态，这一演变可能历时数十亿年。金星的自转周期可能代表一种平衡状态，即太阳引力的潮汐锁定效应（倾向于减慢自转）与厚重大气因太阳加热产生的大气潮汐之间的竞争\(^\text{[158][159]}\)。金星与地球连续两次接近之间的平均间隔为 584 天，这几乎正好等于 5 个金星太阳日（精确值为 5.001444）\(^\text{[160]}\)，但与地球存在自转-公转共振的假说已被否定\(^\text{[161]}\)。

金星没有天然卫星\(^\text{[162]}\)。它有数颗特洛伊小行星：拟卫星 524522 Zoozve\(^\text{[163][164]}\) 以及另外两颗临时特洛伊小行星 2001\,CK32 和 2012\,XE133\(^\text{[165]}\)。17 世纪，乔瓦尼·卡西尼报告称有一颗月亮绕金星运行，并将其命名为 Neith，接下来 200 年中出现了多次观测报告，但大多数后来被确定为附近恒星。加州理工学院 Alex Alemi 和 David Stevenson 在 2006 年对早期太阳系模型的研究显示，金星在数十亿年前一次巨大撞击事件中很可能形成过至少一颗卫星\(^\text{[166]}\)。根据该研究，大约 1000 万年后，另一场撞击事件使金星自转方向逆转，随后的潮汐减速使金星的卫星逐渐向内螺旋并最终与金星相撞\(^\text{[167]}\)。如果后续撞击再次形成过卫星，它们会以同样方式消失。另一种解释金星缺乏卫星的可能性是强烈的太阳潮汐作用会使内侧类地行星周围的大型卫星轨道失稳\(^\text{[162]}\)。

金星轨道空间中存在一条尘埃环云\(^\text{[168]}\)，其可能的来源包括金星追踪的小行星\(^\text{[169]}\)、呈波状迁移的行星际尘埃，或太阳系形成时原始环境盘的遗留物\(^\text{[170]}\)。
\subsubsection{相对于地球的轨道}
\begin{figure}[ht]
\centering
\includegraphics[width=8cm]{./figures/8445fe95a61e302d.png}
\caption{地球位于图示中央，曲线表示金星随时间变化的方向和距离。} \label{fig_Venus_14}
\end{figure}
地球与金星具有近似的 13:8 轨道共振（地球每公转八次，金星公转十三次）\(^\text{[171]}\)。因此，它们彼此接近并在平均 584 天的会合周期中达到下合\(^\text{[5]}\)。从地心视角看，金星相对于地球的轨迹在五个会合周期中绘制出一个五角星形，每个周期旋转 144°。由于该轨迹在视觉上类似花朵，这个金星五角星有时被称为金星之花瓣\(^\text{[172]}\)。

当金星在下合位置位于地球与太阳之间时，它在所有行星中以平均 4100 万\,km（2500 万\,mi）的距离最接近地球\(^\text{[5][note 2][173]}\)。由于地球轨道偏心率正在减小，这一最小距离将在未来数万年中逐渐增大。从公元 1 年到 5383 年，共出现 526 次距离小于 4000 万\,km（2500 万\,mi）的接近；之后约 60,158 年内将不再出现此类接近\(^\text{[174]}\)。

虽然金星接近地球的最近距离最小，但水星更频繁地成为距离地球（以及任何其他行星）最近的行星\(^\text{[175][176]}\)。金星已被用于引力助推的中转点，被认为是前往水星\(^\text{[177][178]}\)、太阳\(^\text{[179]}\)、小行星\(^\text{[180]}\)、火星\(^\text{[181]}\)、木星及更远处\(^\text{[182][183]}\)更快速且经济的航程方式。

在潮汐作用方面，金星对地球施加的潮汐力在月球与太阳之后位列第三，但要弱得多\(^\text{[184]}\)。
\subsection{可观测性}
\begin{figure}[ht]
\centering
\includegraphics[width=8cm]{./figures/149b4a28cac2124e.png}
\caption{金星（图中偏右的位置）从地球观测时在其最大亮度下始终比所有其他行星或恒星更亮。图像顶部可见木星。} \label{fig_Venus_15}
\end{figure}
在肉眼看来，金星呈现为一个白色光点，其最大视星等为 −4.92，比除太阳外的任何行星或恒星都要明亮\(^\text{[185]}\)，即使在凌日过程中最暗时，其视星等也为 −2.98\(^\text{[20]}\)。该行星的平均视星等为 −4.14，标准差为 0.31\(^\text{[20]}\)。在下合前后约一个月的娥眉月相期间，其亮度达到最大。当前方被太阳背光照亮时，金星会变暗至约 −3 等，但确切亮度取决于位相角\(^\text{[186]}\)。金星亮度足以在白昼中被看见\(^\text{[187]}\)，但在太阳位于地平线附近或正在落下时更易观测。作为一颗内行星，它与太阳的角距离始终不超过约 47°\(^\text{[188]}\)。

金星每 584 天在绕太阳运行时“追上”地球\(^\text{[5]}\)。在此过程中，它从日落后可见的“昏星”变为日出前可见的“晨星”。虽然另一颗内行星水星最大延伸角仅为 28°，在曙暮光中往往难以辨认，但金星在最亮时几乎不可能错过。其更大的最大延伸角意味着它在日落后仍能在暗空中被看到。作为夜空中最亮的点状天体，金星经常被误报为“不明飞行物”\(^\text{[189]}\)。

由于金星在下合时会靠近地球，且其轨道相对地球轨道有倾角，它可以出现在比黄道平面北或南超过 8° 的位置，这超过任何其他行星或月球。每八年左右的三月，金星会出现在距离黄道最北的位置，位于双鱼座（例如 2025 年 3 月中旬）；而每八年左右的八月或九月，它会出现在距离黄道最南的位置，位于室女座（例如 2023 年 8 月下旬）。因此在北半球，金星可能位于太阳以北，并在同一天既作为晨星又作为昏星出现。这些南北偏离的时间每年会略微提前，在 30 个周期（240 年）后，这一周期会逐渐被另一个相差三年的周期取代，因此在地球 243 次公转和金星 395 次公转后情况将接近恢复到原来的样子\(^\text{[190]}\)。

金星的月掩现象，即月球遮挡金星，使地球某些地区的观测者无法看到金星，平均每年大约发生两次，有时同一年会发生多次（尽管十分罕见）。
\subsubsection{相位}
\begin{figure}[ht]
\centering
\includegraphics[width=8cm]{./figures/7f0dd42656d60e11.png}
\caption{金星的相位及其视直径的变化} \label{fig_Venus_16}
\end{figure}
在绕太阳运行的过程中，金星在望远镜视野中呈现出类似月相的相位变化。当金星位于太阳的远端（上合）时，呈现为一个较小而“满”的圆面。在它距离太阳达到最大角距时，金星显示为更大的圆面并呈“上弦相”，此时它在夜空中最为明亮。当金星经过地球与太阳之间的近侧时，在望远镜中呈现为更大的细长“娥眉月”。当金星位于地球与太阳之间（下合）时，其视直径最大并呈“朔相”。其大气可通过望远镜观察到，这是由于阳光在其周围发生折射形成的光晕\(^\text{[188]}\)。这些相位在 4 英寸口径的望远镜中即可清晰可见\(^\text{[191]}\)。虽然肉眼能否观察到金星的相位存在争议，但仍有记录显示有人观察到其娥眉形状\(^\text{[192]}\)。
\subsubsection{日间显现}
\begin{figure}[ht]
\centering
\includegraphics[width=8cm]{./figures/b29b2a44adbea1a0.png}
\caption{金星在白昼中常可被肉眼看见，例如在 2015 年 12 月 7 日月掩事件发生前不久所观测到的情形。} \label{fig_Venus_17}
\end{figure}
当金星足够明亮且与太阳具有足够大的角距离时，它在晴朗的白昼天空中可被肉眼轻易看到，尽管大多数人并不知道要去寻找它\(^\text{[193]}\)。天文学家爱德蒙·哈雷在 1716 年计算了其肉眼可见的最大亮度，当时许多伦敦人因在白天看到金星而感到惊慌。法国皇帝拿破仑·波拿巴曾在卢森堡参加一次招待会时目睹过金星的白昼显现\(^\text{[194]}\)。另一次历史性的金星白昼观测发生在 1865 年 3 月 4 日美国总统亚伯拉罕·林肯在华盛顿特区举行就职典礼期间\(^\text{[195]}\)。
\subsubsection{凌日}
\begin{figure}[ht]
\centering
\includegraphics[width=6cm]{./figures/80b71a3a59e0792d.png}
\caption{通过望远镜投射到白色卡片上的 2012 年金星凌日影像} \label{fig_Venus_18}
\end{figure}
金星凌日是指在下合期间，金星出现在太阳前方。由于金星轨道相对于地球轨道略有倾角，因此大多数每隔 1.6 年发生的金星与地球的下合并不会产生金星凌日。因此，只有当下合发生在六月或十二月的若干天内，即金星和地球的轨道与太阳成一直线时，才会出现金星凌日\(^\text{[196]}\)。其结果是金星凌日目前呈现为一个 8 年、105.5 年、8 年和 121.5 年的序列，构成一个 243 年的周期。

金星凌日在历史上具有重要意义，因为它使天文学家能够确定天文单位的大小，从而推算出太阳系的规模。1639 年，杰里迈亚·霍罗克斯通过首次已知的金星凌日观测（继 1631 年人类首次观测到的行星凌日——水星凌日之后）展示了这一点\(^\text{[197]}\)。

自约翰内斯·开普勒在 1621 年计算出其发生时间以来，至今仅有七次金星凌日被观测到。库克船长于 1768 年航行至塔希提岛以记录第三次金星凌日，随后促成了对澳大利亚东海岸的探索\(^\text{[198][199]}\)。

最近的一对金星凌日发生在 2004 年 6 月 8 日和 2012 年 6 月 5–6 日。该凌日可通过许多在线平台进行实时观看，或在适当设备与条件下进行本地观测\(^\text{[200]}\)。再前一对则发生在 1874 年 12 月和 1882 年 12 月。

下一次金星凌日将发生在 2117 年 12 月和 2125 年 12 月\(^\text{[201]}\)。
\subsubsection{灰光现象}
\begin{figure}[ht]
\centering
\includegraphics[width=8cm]{./figures/28c0a4b35d424546.png}
\caption{自 2022 年起，夜辉被认为是灰光现象最可能的成因。在这张可见光与近红外影像中，它最容易被辨认为沿金星边缘的一条亮线\(^\text{[202]}\)。地表及其特征（如图中可见的暗色区域——阿芙洛狄忒高原的奥夫达高地区）在人眼中难以辨认得如此清晰，尽管据报道某些人能够看到，可能是因为他们对地表发光所在光谱具有更高的敏感度\(^\text{[203]}\)。} \label{fig_Venus_19}
\end{figure}
金星观测中的一个长期谜团是所谓的灰光——即在金星呈娥眉月相时，其暗面上出现的微弱亮光。首次关于灰光的观测声称可追溯至 1643 年，但这一亮光的存在从未得到可靠确认。观测者推测，这可能源于金星大气中的电活动，但也可能是错觉，由观察一个明亮的娥眉形天体所引发的生理效应造成\(^\text{[204][140]}\)。灰光现象常在金星出现在傍晚天空时被报告，此时金星的傍晚界线朝向地球。
\subsection{观测与探索历史}
\subsubsection{早期观测}
金星在地球天空中亮度足以裸眼可见，使其成为人类历史上各个文化所认识和识别的经典行星之一，尤其因为它是仅次于太阳和月亮的地球天空中第三亮的天体。由于金星的运动呈现不连续性（因其接近太阳而消失数日，然后在另一侧地平线上重新出现），一些文化没有将金星视为单一天体\(^\text{[205]}\)；相反，他们认为金星在东西地平线上分别是两颗不同的恒星：晨星与昏星\(^\text{[205]}\)。然而，杰姆代特·纳斯尔时期的圆筒印章以及巴比伦第一王朝的阿弥萨杜卡金星泥板都表明，古代苏美尔人已经知道晨星与昏星是同一天体\(^\text{[206][205][207]}\)。
\begin{figure}[ht]
\centering
\includegraphics[width=6cm]{./figures/07505049521d37a0.png}
\caption{已知最古老的金星位置记录，来自阿弥萨杜卡时期的巴比伦金星泥板（公元前 1600 年）} \label{fig_Venus_20}
\end{figure}
在古巴比伦时期，金星被称为 Ninsi'anna，后来被称为 Dilbat\(^\text{[208]}\)。“Ninsi'anna”一名意为“神圣的女士、天空的光辉”，指的是作为最亮可见“星辰”的金星。该名称更早的拼写使用楔形文字 si4（＝SU，意为“呈红色”），其原始含义可能是“天空之红的神圣女士”，指代晨空与昏空的颜色\(^\text{[209]}\)。

中国古代将晨星金星称为“太白”或“启明”，而将昏星金星称为“长庚”\(^\text{[210]}\)。

古希腊人最初认为金星是两颗独立的星：晨星 Phosphorus 与昏星 Hesperus。老普林尼将认识到这两者是同一天体的功劳归于公元前六世纪的毕达哥拉斯\(^\text{[211]}\)，而第欧根尼·拉尔修则认为应归功于巴门尼德（公元前五世纪早期）\(^\text{[212]}\)。尽管古希腊人已认识到金星是同一颗天体，古罗马人仍沿用不同称呼，将其晨星状态称为 Lucifer（字面意为“光之携带者”），将其昏星状态称为 Vesper\(^\text{[213]}\)，两者均为其传统希腊名称的直译。

在二世纪的《天文学大成》中，托勒密提出水星和金星均位于太阳与地球之间。十一世纪波斯天文学家阿维森纳声称曾观测到金星凌日（尽管这一点存在争议）\(^\text{[214]}\)，后来的天文学家将其视为托勒密理论的佐证\(^\text{[215]}\)。十二世纪安达卢西亚天文学家伊本·巴贾赫曾观测到“两个在太阳表面上的黑点”；十三世纪马拉盖天文学家库特布丁·希拉齐认为这些是金星和水星的凌日\(^\text{[216]}\)，但这不可能为真，因为在伊本·巴贾赫的一生中并无金星凌日\(^\text{[note 3]}\)。
\begin{figure}[ht]
\centering
\includegraphics[width=6cm]{./figures/1899ac8f6b035776.png}
\caption{前哥伦布时期的玛雅《德累斯顿抄本》，其中记录并计算了金星的出现时机} \label{fig_Venus_21}
\end{figure}
\subsubsection{金星与近代早期天文学}
\begin{figure}[ht]
\centering
\includegraphics[width=10cm]{./figures/258414620ee3be6f.png}
\caption{1610 年，伽利略通过望远镜观测到金星呈现出不同的相位，尽管它在地球天空中始终靠近太阳（第一幅图）。这一发现证明金星绕太阳运行而非绕地球运行，如哥白尼的日心模型所预言，并推翻了托勒密的地心模型（第二幅图）。} \label{fig_Venus_22}
\end{figure}
当金星于 1610 年 12 月首次被意大利物理学家伽利略用望远镜观测时，他发现金星呈现如月亮般的相位变化，从娥眉月到凸月再到满相，反之亦然。当金星在天空中远离太阳的位置时，它呈现半月相；而当它靠近太阳时，则呈现娥眉相或满相。只有当金星绕太阳运行时，这种现象才可能出现。伽利略在 1613 年的《太阳黑子书信》中报告了这一发现，这成为最早明确反驳托勒密地心体系（即太阳系同心且以地球为中心）的观测之一\(^\text{[219][220]}\)。

1631 年的金星凌日虽然未被记录下来，但这是首次被成功预测的凌日，由约翰内斯·开普勒计算并于 1629 年发表。接下来的 1639 年金星凌日由杰里迈亚·霍罗克斯精确预测，并由他与朋友威廉·克拉布特里各自在家中于 1639 年 12 月 4 日（当时使用的儒略历为 11 月 24 日）同时观测到\(^\text{[221]}\)。
\begin{figure}[ht]
\centering
\includegraphics[width=8cm]{./figures/802be5669f872ee3.png}
\caption{描绘杰里迈亚·霍罗克斯观测 1639 年金星凌日的二十世纪绘画} \label{fig_Venus_23}
\end{figure}
金星的大气由俄罗斯博学家米哈伊尔·罗蒙诺索夫于 1761 年发现\(^\text{[222][223]}\)。1790 年，德国天文学家约翰·施雷特尔观测了金星的大气。施雷特尔发现，当金星呈现细娥眉相时，其两端的尖角延伸超过 180°。他正确地推断这源于致密大气中阳光的散射。之后，美国天文学家切斯特·史密斯·莱曼在金星处于下合时观察到其暗面周围的完整光环，为其具有大气提供了进一步证据\(^\text{[224]}\)。大气的存在使确定金星自转周期的工作变得复杂，意大利出生的天文学家乔瓦尼·卡西尼和施雷特尔等观测者根据金星表观表面标记的运动错误地估计其自转周期约为 24 小时\(^\text{[225]}\)。
\begin{figure}[ht]
\centering
\includegraphics[width=6cm]{./figures/d02c63b404bf6fdf.png}
\caption{1769 年金星凌日期间记录到的“黑滴效应”} \label{fig_Venus_24}
\end{figure}
\subsubsection{二十世纪早期的进展}
在二十世纪之前，人们对金星几乎没有新的认识。其看似毫无特征的圆盘状外观未能透露其地表状况，直到光谱学和紫外观测的发展才揭示了更多秘密。

最早的紫外观测于 1920 年代进行，当时弗兰克·E·罗斯发现紫外照片中呈现出大量在可见光和红外波段中不存在的细节。他认为这是由于一层致密的黄色低层大气，上方覆盖着高层卷云所致\(^\text{[226]}\)。

人们注意到金星的圆盘缺乏可辨识的扁率，暗示其自转较慢，一些天文学家据此认为金星像当时被认为潮汐锁定的水星一样被潮汐锁定；但另一些研究者探测到金星夜侧存在大量热量，这又暗示其自转速度很快（当时人们尚未意识到其极高的表面温度），使问题更加复杂\(^\text{[227]}\)。后来在 1950 年代的研究表明，金星的自转为逆行。
\subsubsection{首次金星探测任务}
1961 年，人类首次尝试进行行星际空间飞行，即苏联“金星”计划中的无人探测器“金星 1 号”飞向金星，但在途中失去联系\(^\text{[228]}\)。

第一次成功的行星际任务，同样是前往金星，是美国“水手”计划的“水手 2 号”，于 1962 年 12 月 14 日以距金星表面 34,833\,km（21,644\,mi）的距离飞掠，并获取了金星大气的数据\(^\text{[229][230]}\)。

此外，最早的金星雷达观测在 1960 年代实施，并首次测得了其自转周期，其数值接近真实值\(^\text{[231]}\)。

1966 年发射的“金星 3 号”成为人类首个抵达并撞击月球以外其他天体的探测器和着陆器，但因坠毁于金星表面而未能返回数据。1967 年发射的“金星 4 号”成功在金星大气中实施科学实验后着陆。金星 4 号显示，金星表面温度几乎达 500\,°C（932\,°F），比水手 2 号计算的更高；其测得大气成分为 95\% 的二氧化碳（$CO_{2}$），并发现金星大气密度远超设计者的预期\(^\text{[232][233]}\)。

作为早期太空合作的范例，金星 4 号的数据与 1967 年的水手 5 号数据相结合，并由苏美联合科学团队在接下来一年的一系列专题会议中共同分析\(^\text{[234]}\)。

1970 年 12 月 15 日，“金星 7 号”成为首个在另一颗行星上软着陆并将数据传回地球的航天器\(^\text{[235]}\)。

1974 年，“水手 10 号”飞掠金星，以借助金星引力改变航向飞往水星，并拍摄了金星云层的紫外影像，揭示了金星大气中极高的风速。这是历史上首次使用的行星际引力助推技术，后来被众多探测器沿用。

1970 年代的雷达观测首次揭示了金星表面的地形细节。研究人员利用阿雷西博天文台 300\,m（1000\,ft）射电望远镜向金星发射无线电波，回波显示出两个反射率极高的区域，被命名为 Alpha 区和 Beta 区。观测还发现一个亮区，被归因于山地，并被命名为麦克斯韦山脉（Maxwell Montes）\(^\text{[236]}\)。这三处特征是目前金星上仅有的以非女性名称命名的地表特征\(^\text{[50]}\)。
\begin{figure}[ht]
\centering
\includegraphics[width=8cm]{./figures/9b4b2711eb58049a.png}
\caption{金星表面的首次影像及首次清晰的 180 度全景图，也是除地球外其他行星的首张此类全景图（1975 年，苏联“金星 9 号”着陆器）。黑白图像中呈现荒芜、黑色、板岩状的岩石，背景是一片平坦的天空。地面与探测器为画面主体。} \label{fig_Venus_26}
\end{figure}
1975 年，苏联的“金星 9 号”和“金星 10 号”着陆器传回了来自金星表面的首批图像，这些图像为黑白影像。美国宇航局随后通过“先驱者金星”计划获取了更多数据，该计划由两个独立任务组成\(^\text{[237]}\)：“先驱者金星多探测器”与“先驱者金星轨道器”，二者在 1978 年至 1992 年间绕金星运行\(^\text{[238]}\)。1982 年，苏联“金星 13 号”和“金星 14 号”着陆器首次获取了金星表面的单色彩色滤光片图像。在 1983 至 1984 年间，“金星 15 号”和“金星 16 号”在金星轨道运行，对其 25\% 的地表（从北极至北纬 30°）开展了详细测绘，之后苏联的“金星”探测计划宣告结束\(^\text{[239]}\)。
\begin{figure}[ht]
\centering
\includegraphics[width=6cm]{./figures/2cf993976ba79b80.png}
\caption{史密森学会乌德瓦－哈兹中心展出的“织女号”气球探测器} \label{fig_Venus_27}
\end{figure}
1985 年，苏联的“织女号”计划（Vega 1 与 Vega 2）携带了最后一批进入式探测器，并首次携带了地球以外的空气机器人（aerobots），通过充气气球的方式首次实现了地外大气中的飞行。

1990 至 1994 年间，“麦哲伦号”在金星轨道运行直至轨道衰减，并对金星表面进行了测绘。此外，“伽利略号”（1990）\(^\text{[240]}\)与“卡西尼–惠更斯号”（1998/1999）等探测器在飞往其他目的地途中也曾飞掠金星。
\subsubsection{重新开展的探索}
2006 年 4 月，欧洲航天局（ESA）的首次金星专门探测任务“金星快车”进入金星轨道。“金星快车”为金星大气带来了前所未有的观测成果。ESA 于 2014 年 12 月结束该任务，并于 2015 年 1 月使其再入\(^\text{[241]}\)。同年及次年，“信使号”（MESSENGER）在前往其他目的地途中也进行了金星飞掠。

2010 年，首艘成功的行星际太阳帆飞船 IKAROS 前往金星并实施了飞掠。

2015 至 2024 年间，日本的“曙光号”（Akatsuki）在金星轨道上开展探测，同时“贝皮科伦布号”（BepiColombo）在 2020/2021 年进行了飞掠。
\begin{figure}[ht]
\centering
\includegraphics[width=8cm]{./figures/51a8621358abd9f4.png}
\caption{于 2021 年拍摄了这段金星夜侧的可见光影像，透过云层显示出其炽热且微弱发光的表面，以及作为大片暗区的阿芙洛狄忒高地；而在白昼侧，由于云层被照亮，此类观测无法进行\(^\text{[242][243]}\)。} \label{fig_Venus_28}
\end{figure}
\subsubsection{正在执行与计划中的任务}
\begin{figure}[ht]
\centering
\includegraphics[width=8cm]{./figures/25fe75bbacd989dd.png}
\caption{金星全球地形图，并标注所有探测器着陆点（红色：返回图像；黑点：采集样本并进行现场分析）} \label{fig_Venus_29}
\end{figure}
截至 2025 年，金星附近没有正在运行的探测器，帕克太阳探测器计划在 2030 年之前多次返回金星进行飞掠。

多项探测器正在研制中，同时还有多项任务仍处于早期概念阶段。下一项已计划的金星任务是“金星生命探测器”（Venus Life Finder），预计最早于 2026 年夏季发射。

印度空间研究组织（ISRO）正在推进金星轨道器任务（Venus Orbiter Mission），计划于 2028 年发射。阿联酋的 MBR 探索者（MBR Explorer）小行星任务将对金星进行一次飞掠。美国宇航局（NASA）已批准两项金星任务——VERITAS 与 DAVINCI——计划最早于 2031 年发射。欧洲航天局（ESA）计划于 2031 年发射 EnVision。

\textbf{目标}

金星被确定为未来研究中的重要对象，用于理解：
\begin{itemize}
\item 太阳系与地球的起源，以及类似我们这样的行星系统在宇宙中是普遍还是稀有；
\item 行星天体如何从原始状态演化至今日多样的形态；
\item 导致适宜生存环境与生命出现的条件如何发展[244]。
\end{itemize}
\subsubsection{载人任务概念}
自 1960 年代起，金星便被视作火星载人探测的潜在中转点，即采用“冲日任务”，而非直接利用金星引力辅助的“合日任务”。研究显示，这类方案使前往火星的任务更快、更安全，并拥有更好的返程或中止窗口，且辐射暴露不高于或甚至低于直接前往火星的任务\(^\text{[245][246]}\)。

\textbf{金星大气层中的可能居住}
\begin{figure}[ht]
\centering
\includegraphics[width=8cm]{./figures/2019ebdd05da8040.png}
\caption{美国宇航局“金星高空作业构想”（HAVOC）载人浮空前哨基地的艺术示意图} \label{fig_Venus_30}
\end{figure}
尽管金星表面的环境极端恶劣，但在距地表 50\,km 的高度，其大气压、温度以及所受的太阳与宇宙辐射均与地球表面相似（“温和条件”）\(^\text{[247][114][113][183]}\)。然而，若要在金星大气中建立任何形式的人类存在，面临的众多工程挑战之一是其大气中具有强腐蚀性的高浓度硫酸\(^\text{[248]}\)。作为替代在火星等行星表面生活的流行设想，人们提出了用于载人探索甚至永久“浮空城市”的金星大气层浮空器方案\(^\text{[248][249][250][251][252]}\)。美国宇航局的“金星高空作业构想”（HAVOC）则是一项用于研究载人浮空器设计的训练性概念。
\subsection{生命的可能性}
自 20 世纪 60 年代初以来，关于金星表面存在生命的推测显著减少，因为人们逐渐认识到金星表面的环境相比地球极端得多。金星的极端温度与大气压使目前已知的以水为基础的生命形式难以存在。

一些科学家推测，在金星较高、较冷且酸性的上层大气中，可能存在嗜热嗜酸的极端微生物\(^\text{[253][254][255]}\)。这类推测可追溯到 1967 年，当时卡尔·萨根和哈罗德·J·莫罗维茨在《自然》杂志发文指出，金星云层中探测到的微小颗粒或许是类似地球细菌的生物（其尺度相似）：

“虽然金星地表条件使生命假说不太可信，但金星云层则是完全不同的故事。正如数年前有人指出的那样，水、二氧化碳和阳光——光合作用所需的要素——在云层附近极为丰富。”\(^\text{[256]}\)

2019 年 8 月，由李延珠（Yeon Joo Lee）领导的天文学团队报告称，金星大气中“未知吸收体”引起了长期的吸收率与反照率变化，这些吸收体可能是化学物质，也可能是在高空存在的大型微生物群落，它们会影响金星的气候\(^\text{[137]}\)。其光吸收曲线与地球云层中的微生物几乎完全一致。其他研究也得出了类似结论\(^\text{[257]}\)。

2020 年 9 月，由卡迪夫大学的简·格里夫斯（Jane Greaves）领导的天文学团队宣布，他们可能在金星上层云层中探测到磷化氢，这是一种已知不会由金星地表或大气中的化学过程产生的气体\(^\text{[258][116][115][259][260]}\)。一种可能的解释是生命活动\(^\text{[261]}\)。该磷化氢在距地表至少 48\,km（30\,mi）高度、主要位于中纬度地区被探测到，而在两极则没有发现。这一发现促使 NASA 管理员吉姆·布里登斯廷公开呼吁重新关注金星研究，称磷化氢的发现是“在人类寻找地外生命道路上最重要的发展之一”\(^\text{[262][263]}\)。

随后对用于识别金星大气中磷化氢的数据处理方式的分析提出了质疑：用于拟合的 12 阶多项式可能放大了噪声并生成了假信号（参见龙格现象）。在其他电磁波段中对金星大气的观测未能检出预期的磷化氢吸收线\(^\text{[264]}\)。到 2020 年 10 月下旬，重新分析并正确扣除背景后，已不再显示具有统计显著性的磷化氢信号\(^\text{[265][266][267]}\)。

格里夫斯团队的部分成员正在参与麻省理工学院的计划，准备与 Rocket Lab 火箭公司合作发射首个私人行星际探测器“金星生命探测器”（Venus Life Finder），通过探测器进入金星大气以寻找有机物\(^\text{[269]}\)。
\subsubsection{行星保护}
由于金星表面环境极其严酷，金星被归类为行星保护第二类（category two），即第二低级别\(^\text{[269]}\)。这意味着航天器携带的星际污染几乎不可能影响对其的科学研究。

然而，随着金星可能存在生物标志物的发现，对至少部分金星大气层的分类提出了质疑。但由于这些层尚未被确认为能够支持生命，因此不建议调整其行星保护分类\(^\text{[270]}\)。
\subsubsection{人类可居性}
尽管金星表面对人类极不友好，但在距地表 50\,km 的高度，其大气条件不仅被认为可能适宜本地生命，也可能适宜人类生存，比太阳系除地球以外的任何地方都更具优势。从大气压、重力、温度到辐射，除了化学成分外，这里的条件几乎与地球表面相同。基于这一前景，有人提出可在此高度建造人类的浮空居住设施，以实现人类前往金星的可能性\(^\text{[271]}\)。
\subsection{文化中的金星}
金星是夜空中最显著的天体之一，在众多文化的神话、占星学与文学作品中都被赋予了特殊的重要地位。
\begin{figure}[ht]
\centering
\includegraphics[width=8cm]{./figures/a5aa0ac67bd47cec.png}
\caption{} \label{fig_Venus_31}
\end{figure}
多首赞美诗歌颂伊南娜作为金星女神的角色\(^\text{[205][272][273]}\)。神学教授杰弗里·库利（Jeffrey Cooley）认为，在许多神话中，伊南娜的行动可能与金星在天空中的运行相对应\(^\text{[205]}\)。金星不连续的运动与伊南娜的神话形象及其双重性相关\(^\text{[205]}\)。在《伊南娜降入冥界》中，伊南娜能够降入冥界又重返天界，这是其他神祇无法做到的。金星似乎也有类似的“降落”，它在西方落下，又在东方升起\(^\text{[205]}\)。一首引言性赞美诗描述伊南娜离开天界前往库尔（Kur），被认为是山地，其意象与伊南娜在西方的升落相对应\(^\text{[205]}\)。《伊南娜与舒卡勒图达》《伊南娜降入冥界》中都出现了与金星运行相平行的描绘\(^\text{[205]}\)。《伊南娜与舒卡勒图达》中，舒卡勒图达被描述为在天穹中寻找伊南娜，可能指在东方与西方地平线上寻找\(^\text{[274]}\)。在同一神话中，伊南娜在寻找攻击者时的多次移动也与金星在天空中的运行一致\(^\text{[205]}\)。

通过美索不达米亚文化的影响，古埃及人与希腊人可能在公元前二千纪（或最晚在更后的晚期）便知道晨星与昏星为同一天体\(^\text{[275][276]}\)。埃及人称晨星为 Tioumoutiri，昏星为 Ouaiti\(^\text{[277]}\)。他们起初将金星描绘为凤凰或苍鹭（见 Bennu）\(^\text{[275]}\)，称其为“穿越者”或“带十字之星”\(^\text{[275]}\)，并与奥西里斯相联系；后又将其描绘为双头形象（人首或隼首），并与荷鲁斯相关联\(^\text{[276]}\)，荷鲁斯是伊西斯之子（在更后的希腊化时期，伊西斯与哈索尔被认同为阿佛洛狄忒）。希腊人将晨星称为 Phōsphoros（光明使者，磷元素名“phosphorus”即源于此；或 Ēōsphoros，意为“黎明带来者”），将昏星称为 Hesperos（意为“西方者”）\(^\text{[278]}\)，二者皆为黎明女神 Eos 之子，因此是阿佛洛狄忒的外孙。尽管到罗马时代二者已被认作同一天体并统称为“维纳斯之星”，但传统的两个希腊名称仍然沿用，并通常以拉丁语译为 Lucifer 与 Vesper\(^\text{[278][279]}\)。

古典诗人如荷马、萨福、奥维德与维吉尔都描写过金星及其光辉\(^\text{[280]}\)。诗人威廉·布莱克、罗伯特·弗罗斯特、莱蒂希娅·伊丽莎白·兰登、阿尔弗雷德·丁尼生与威廉·华兹华斯也都为金星写过颂歌\(^\text{[281]}\)。作曲家霍尔斯特在其《行星组曲》中将金星作为第二乐章。

在印度，金星的梵文名称为 Shukra Graha，意为“圣者 Shukra 之星”，Shukra 是一位强大的圣人。在吠陀占星学中\(^\text{[282]}\)，Shukra 的梵文含义为“清澈、纯净”或“明亮、清晰”。作为九曜之一，它被认为影响财富、享乐与繁衍；它是 Bhrgu 之子，Daityas 的导师，也是 Asuras 的古鲁\(^\text{[283]}\)。

英语名称“Venus”源自古罗马人。罗马人以爱神维纳斯为该行星命名，而维纳斯的原型是希腊的爱神阿佛洛狄忒\(^\text{[284]}\)，而阿佛洛狄忒又源自苏美尔宗教中的伊南娜（阿卡德宗教中称为伊什塔尔），她们皆与金星相关\(^\text{[273][272]}\)。金星与这些女神对应的星期为星期五（Friday），以日耳曼女神 Frigg 命名，而 Frigg 也与罗马的维纳斯相关。

中文中，金星称为“金星”，意为“金之星”；在中国传统五行中，金星属金。这一传统在现代中、日、韩、越文化中共享，各语言中金星均字面意为“金星”（金星）\(^\text{[285][286][287][288]}\)。

玛雅人认为金星仅次于太阳和月亮，是最重要的天体。他们称金星为 Chac ek\(^\text{[289]}\)或 Noh Ek'（意为“大星”）\(^\text{[290]}\)。金星的周期对其历法十分重要，并在他们的一些典籍中记载，例如《墨西哥玛雅手抄本》和《德累斯顿抄本》。智利国旗（孤星旗）所描绘的“孤星”即象征金星。
\subsubsection{现代文化}
\begin{figure}[ht]
\centering
\includegraphics[width=8cm]{./figures/47a31ae5a698ebd5.png}
\caption{在文森特·梵高 1889 年的绘画《星夜》中，金星被描绘在画面中那棵巨大柏树的右侧\(^\text{[291][292]}\)。} \label{fig_Venus_32}
\end{figure}
无法穿透的金星云层使科幻作家得以自由想象其表面状况；尤其在早期观测显示金星不仅与地球大小相近，还拥有稠密大气之后，他们更是大胆发挥。由于金星比地球更靠近太阳，它常被描绘为更温暖，但仍适宜人类居住\(^\text{[293]}\)。这一文学类型在 1930 至 1950 年代达到巅峰，当时科学已揭示金星的一些特征，但尚不了解其表面的严酷现实。首次金星探测任务的发现证明真实情况截然不同，从而终结了这一体裁\(^\text{[294]}\)。随着对金星科学认知的进步，科幻作家也试图同步发展，尤其通过描绘人类试图地球化（金星改造）金星的设想\(^\text{[295]}\)。
\subsubsection{符号}
圆圈下带一小十字的符号即所谓的金星符号，因作为金星的天文符号而得名。该符号源于古希腊，更广泛地代表女性特质，并被生物学采用为女性性别符号\(^\text{[296][297][298]}\)，与用于男性的火星符号相对应，有时也与用作雌雄同体现象的水星符号并列。这种将金星与火星赋以性别含义的做法，也被用于异性恋规范化的配对语境中，用以将女性与男性刻板地描述为来自完全不同的“行星”。这种理解在 1992 年出版的《男人来自火星，女人来自金星》一书中被进一步普及\(^\text{[299]}\)。

金星符号也曾用于西方炼金术中，代表元素铜（如同水星符号亦代表元素汞）\(^\text{[297][298]}\)。由于抛光铜自古以来一直用于制作镜子，金星符号有时被称为“金星之镜”，象征爱神的镜子，尽管这一解释后来被认为不太可能是其真正来源\(^\text{[297][298]}\)。

除金星符号外，还有许多其他符号与金星相关，其中常见的还有新月，或尤以星形符号著称，如伊什塔尔之星\(^\text{[300]}\)。
\subsection{另见}
\begin{itemize}
\item 金星纲要  
\item 太阳系行星的物理性质  
\item 金星带 
\end{itemize} 
\subsection{注释 } 
\begin{enumerate}
\item 在新闻稿与科学出版物中被误写为 “Ganiki Chasma”\(^\text{[66]}\)。  
\item 必须明确“接近度”一词的含义。在天文学文献中，“最近行星”（closest planets）常指两颗行星在其轨道上彼此最接近的程度，即两者轨道之间的最小距离。然而，这并不意味着这两颗行星在时间尺度上总体上彼此最近。由于水星比金星更靠近太阳，水星在更长时间内与地球保持相对接近；因此，从时间平均的角度说，可以认为“水星是与地球平均意义上最近的行星”。然而，在采用这种时间平均定义时，水星实际上是整个太阳系中与所有其他行星平均距离最近的行星。因此，这种基于平均距离的“接近”定义似乎并不特别有用。BBC Radio 4 节目《More or Less》的某一集对这些不同的“接近”概念进行了相当清晰的解释\(^\text{[173]}\)。  
\item 若干中世纪伊斯兰天文学家声称的凌日观测后来被证明是太阳黑子[217]。阿维森纳并未记录他的观测日期。在他的一生中确实发生过一次金星凌日，即 1032 年 5 月 24 日，但尚不确定当时从他的所在地是否能够看到这一现象\(^\text{[218]}\)。
\end{enumerate}
\subsection{参考文献}
\begin{enumerate}
\item Ricardo, Nunes. “Ricardo Nunes Astronomy Page”。[www.astrosurf.com。检索于](http://www.astrosurf.com。检索于) 2025 年 11 月 19 日。
\item “Venusian”。Lexico 英国英语词典。牛津大学出版社。2020 年 3 月 23 日存档。
\item “Venusian”。Merriam-Webster.com 词典。Merriam-Webster。
\item “Cytherean”。《牛津英语词典》（在线版）。牛津大学出版社。（需订阅或机构访问权限。）
\item “Venerean, Venerian”。《牛津英语词典》（在线版）。牛津大学出版社。（需订阅或机构访问权限。）
\item Williams, David R.（2020 年 11 月 25 日）。“Venus Fact Sheet”。美国宇航局戈达德太空飞行中心。2018 年 5 月 11 日存档。2021 年 4 月 15 日检索。
\item Yeomans, Donald K. “Horizons Web-Interface for Venus (Major Body=2)”。JPL Horizons 在线星历系统。2023 年 8 月 18 日存档。2010 年 11 月 30 日检索。——选择“Ephemeris Type: Orbital Elements”“Time Span: 2000-01-01 12:00 to 2000-01-02”。（“Target Body: Venus”“Center: Sun”保持默认。）结果为 J2000 历元的瞬时摄动轨道元素。
\item Simon, J.L.; Bretagnon, P.; Chapront, J.; Chapront-Touzé, M.; Francou, G.; Laskar, J.（1994 年 2 月）。“月球与行星的岁差公式与平均轨道要素的数值表达式”。《天文学与天体物理学》282（2）：663–683。Bibcode:1994A&A...282..663S。
\item Souami, D.; Souchay, J.（2012 年 7 月）。“太阳系的不变量面”。《天文学与天体物理学》543：11。Bibcode:2012A&A...543A.133S。doi:10.1051/0004-6361/201219011。A133。
\item Seidelmann, P. Kenneth; Archinal, Brent A.; A'Hearn, Michael F.; 等（2007）。“IAU/IAG 工作组关于行星和卫星制图坐标和自转参量（2006）的报告”。《天体力学与动力天文学》98（3）：155–180。Bibcode:2007CeMDA..98..155S。doi:10.1007/s10569-007-9072-y。
\item Konopliv, A. S.; Banerdt, W. B.; Sjogren, W. L.（1999 年 5 月）。“金星引力：180 阶度模型”（PDF）。《伊卡洛斯》139（1）：3–18。Bibcode:1999Icar..139....3K。CiteSeerX 10.1.1.524.5176。doi:10.1006/icar.1999.6086。2010 年 5 月 26 日从原 PDF 存档。
\item “行星与冥王星：物理特征”。NASA。2008 年 11 月 5 日。2006 年 9 月 7 日存档。2015 年 8 月 26 日检索。
\item “行星事实”。行星学会（The Planetary Society）。2012 年 5 月 11 日存档。2016 年 1 月 20 日检索。
\item Margot, Jean-Luc; Campbell, Donald B.; Giorgini, Jon D.; 等（2021 年 4 月 29 日）。“金星的自旋状态与转动惯量”。《自然·天文学》（Nature Astronomy）。5（7）：676–683。arXiv:2103.01504。Bibcode:2021NatAs...5..676M。doi:10.1038/s41550-021-01339-7。S2CID 232092194。
\item “IAU/IAG 行星与卫星制图坐标和自转要素工作组报告”。国际天文学联合会。2000 年。2020 年 5 月 12 日存档。2007 年 4 月 12 日检索。
\item Archinal, B. A.; Acton, C. H.; A'Hearn, M. F.; Conrad, A.; Consolmagno, G. J.; Duxbury, T.; Hestroffer, D.; Hilton, J. L.; Kirk, R. L.; Klioner, S. A.; McCarthy, D.; Meech, K.; Oberst, J.; Ping, J.; Seidelmann, P. K.（2018）。“IAU 制图坐标与自转要素工作组报告：2015”。《天体力学与动力天文学》。130（3）：22。Bibcode:2018CeMDA.130...22A。doi:10.1007/s10569-017-9805-5。ISSN 0923-2958。
\item Mallama, Anthony; Krobusek, Bruce; Pavlov, Hristo（2017）。“行星的宽波段星等与反照率的综合研究，并应用于系外行星和第九行星假说”。《伊卡洛斯》（Icarus）。282：19–33。arXiv:1609.05048。Bibcode:2017Icar..282...19M。doi:10.1016/j.icarus.2016.09.023。S2CID 119307693。
\item Haus, R.; Kappel, D.; Arnoldb, G.（2016 年 7 月）。“基于改进的金星中低层大气模型的辐射能量平衡”（PDF）。《伊卡洛斯》（Icarus）。272：178–205。Bibcode:2016Icar..272..178H。doi:10.1016/j.icarus.2016.02.048。2017 年 9 月 22 日 PDF 存档。2019 年 6 月 25 日检索。
\item “大气与行星温度”。美国化学学会。2013 年 7 月 18 日。2023 年 1 月 27 日存档。2023 年 1 月 3 日检索。Herbst, K.; Banjac, S.; Atri D.; Nordheim, T. A（2020 年 1 月 1 日）。“在宜居性背景下重新评估宇宙射线导致的金星辐射剂量”。《天文学与天体物理学》（Astronomy & Astrophysics）。633。图 6。arXiv:1911.12788。Bibcode:2020A&A...633A..15H。doi:10.1051/0004-6361/201936968。ISSN 0004-6361。S2CID 208513344。2021 年 12 月 5 日存档。2021 年 11 月 20 日检索。
\item Mallama, Anthony; Hilton, James L.（2018 年 10 月）。“为《天文年鉴》计算行星视星等”。《天文学与计算》（Astronomy and Computing）。25：10–24。arXiv:1808.01973。Bibcode:2018A&C....25...10M。doi:10.1016/j.ascom.2018.08.002。S2CID 69912809。
\item “百科全书——最亮的天体”。IMCCE。2023 年 7 月 24 日存档。2023 年 5 月 29 日检索。
\item Lawrence, Pete（2005）。《寻找金星之影》。Digitalsky.org.uk。2012 年 6 月 11 日存档。2012 年 6 月 13 日检索。
\item Walker, John。《在大白天观测金星》。Fourmilab Switzerland。2017 年 3 月 29 日存档。2017 年 4 月 19 日检索。
\item Andrews, Roy George（2024 年 5 月 27 日）。《金星上的熔岩河揭示更为活跃的火山行星——新软件让科学家重新分析旧雷达影像，提供了火山仍在重塑这颗地狱般行星的最强证据之一》。《纽约时报》。2024 年 5 月 28 日存档。2024 年 5 月 28 日检索。
\item Sulcanese, Davide; Mitri, Giuseppe; Mastrogiuseppe, Marco（2024 年 5 月 27 日）。《麦哲伦雷达揭示金星正在进行的火山活动证据》。《自然·天文学》（Nature Astronomy）。8（8）：973–982。Bibcode:2024NatAs...8..973S。doi:10.1038/s41550-024-02272-1。ISSN 2397-3366。2024 年 5 月 28 日存档。2024 年 5 月 28 日检索。
\item Chang, Kenneth（2023 年 10 月 26 日）。《数十亿年前，金星或曾拥有关键的类地特征——新研究提出，这颗地狱般的第二行星可能曾有板块构造，使其更适宜生命》。《纽约时报》。2023 年 10 月 26 日存档。2023 年 10 月 27 日检索。
\item Weller, Matthew B.; 等（2023 年 10 月 26 日）。《金星大气中的氮由古代板块构造解释》。《自然·天文学》（Nature Astronomy）。7（12）：1436–1444。Bibcode:2023NatAs...7.1436W。doi:10.1038/s41550-023-02102-w。S2CID 264530764。2023 年 10 月 27 日存档。2023 年 10 月 27 日检索。
\item Jakosky, Bruce M.（1999）。《类地行星的大气》。载于 Beatty, J. Kelly；Petersen, Carolyn Collins；Chaikin, Andrew（编）。《新太阳系》（第 4 版）。波士顿：Sky Publishing。第 175–200 页。ISBN 978-0-933346-86-4。OCLC 39464951。
\item Hashimoto, George L.; Roos-Serote, Maarten; Sugita, Seiji; Gilmore, Martha S.; Kamp, Lucas W.; Carlson, Robert W.; Baines, Kevin H.（2008 年 12 月 31 日）。《伽利略近红外成像光谱仪数据暗示金星具有长英质高地地壳》。《地球物理研究杂志：行星》（Journal of Geophysical Research: Planets）。113（E5），2008JE003134。《地球与空间科学进展》。Bibcode:2008JGRE..113.0B24H。doi:10.1029/2008JE003134。S2CID 45474562。
\item Shiga, David（2007 年 10 月 10 日）。《金星的古老海洋是否孕育了生命？》。《新科学家》。2009 年 3 月 24 日存档。2017 年 9 月 17 日检索。
\item Lopes, Rosaly M. C.; Gregg, Tracy K. P.（2004）。《火山世界：探索太阳系火山》。施普林格出版社。第 61 页。ISBN 978-3-540-00431-8。
\item Squyres, Steven W.（2016）。“金星”。《大英百科全书在线》。2014 年 4 月 28 日存档。2016 年 1 月 7 日检索。
\item Darling, David。《金星》。《科学百科全书》。英国邓迪。2021 年 10 月 31 日存档。2022 年 3 月 24 日检索。
\item Lammer, Helmut; Brasser, Ramon; Johansen, Anders; Scherf, Manuel; Leitzinger, Martin（2021 年 2 月）。《金星、地球与火星的形成：由同位素约束》。《空间科学评论》（Space Science Reviews）。217（1）：7。arXiv:2102.06173。doi:10.1007/s11214-020-00778-4。ISSN 0038-6308。
\item “第 6 章：迈向巨型化——NASA 科学”。2018 年 12 月 11 日。2025 年 9 月 7 日检索。
\item Koren, Marina（2022 年 2 月 3 日）。《我们的太阳系以真实色彩呈现时竟如此不同凡响》。《大西洋月刊》。2025 年 6 月 24 日检索。
\item Mueller, Nils（2014）。《金星的表面与内部》。见 Tilman Spohn；Doris Breuer；T. V. Johnson（编），《太阳系百科全书》（第 3 版）。牛津：爱思唯尔科学与技术出版社。ISBN 978-0-12-415845-0。2021 年 9 月 29 日存档。2016 年 1 月 12 日检索。
\item Esposito, Larry W.（1984 年 3 月 9 日）。《二氧化硫：间歇性喷注显示金星存在活跃火山作用的证据》。《科学》（Science）。223（4640）：1072–1074。Bibcode:1984Sci...223.1072E。doi:10.1126/science.223.4640.1072。PMID 17830154。S2CID 12832924。
\item Bullock, Mark A.; Grinspoon, David H.（2001 年 3 月）。《金星气候的近期演化》（PDF）。《伊卡洛斯》（Icarus）。150（1）：19–37。Bibcode:2001Icar..150...19B。CiteSeerX 10.1.1.22.6440。doi:10.1006/icar.2000.6570。2003 年 10 月 23 日存档（PDF）。
\item Basilevsky, Alexander T.; Head, James W. III（1995）。《金星的全球地层学：对三十六个随机测试区域的分析》。《地球、月球与行星》（Earth, Moon, and Planets）。66（3）：285–336。Bibcode:1995EM&P...66..285B。doi:10.1007/BF00579467。S2CID 21736261。

\item Jones, Tom；Stofan, Ellen（2008）。《行星学：解锁太阳系的奥秘》。国家地理学会。第 74 页。ISBN 978-1-4262-0121-9。2017 年 7 月 16 日存档。2017 年 4 月 20 日检索。
\item Kaufmann, W. J.（1994）。《宇宙》。纽约：W. H. Freeman。第 204 页。ISBN 978-0-7167-2379-0。
\item 《国家地理》（2024）。《金星具有活跃火山生命》。
\item Nimmo, F.; McKenzie, D.（1998）。《金星的火山作用与构造》。《年度地球与行星科学评论》。26（1）：23–53。Bibcode:1998AREPS..26...23N。doi:10.1146/annurev.earth.26.1.23。S2CID 862354。
\item Strom, Robert G.; Schaber, Gerald G.; Dawson, Douglas D.（1994 年 5 月 25 日）。《金星的全球再表面化》。《地球物理研究杂志》。99（E5）：10899–10926。Bibcode:1994JGR....9910899S。doi:10.1029/94JE00388。2020 年 9 月 16 日存档。2019 年 6 月 25 日检索。
\item Frankel, Charles（1996）。《太阳系的火山》。剑桥大学出版社。ISBN 978-0-521-47770-3。2023 年 1 月 30 日检索。
\item Mitchell, Don P.（1975 年 10 月 20 日）。《苏联金星图像》。Don P. Mitchell 网站。2025 年 6 月 2 日检索。
\item Moroz, V. I.; Golovin, Yu. M.; Ekonomov, A. P.; Moshkin, B. E.; Parfent'ev, N. A.; San'ko, N. F.（1980）。《金星白天天空的光谱》。《自然》（Nature）。284（5753）：243–244。Bibcode:1980Natur.284..243M。doi:10.1038/284243a0。ISSN 0028-0836。2025 年 6 月 2 日检索。
\item Batson, R. M.; Russell, J. F.（1991 年 3 月 18–22 日）。《为金星新发现的地貌命名》（PDF）。《第 22 届月球与行星科学会议论文集》。美国德克萨斯州休斯敦。第 65 页。Bibcode:1991pggp.rept..490B。2011 年 5 月 13 日存档（PDF）。2009 年 7 月 12 日检索。


\end{enumerate}